%%% BEGIN fig:dis-oneloop
\begin{figure}[htbp]
\centering
% photons, photon vertices, cut and external labels common to the cut graphs (a)-(d)
\newcommand{\toneDIS}{%
  \draw[photon] (-1.65,2.6) -- (-1.15,1.8);
  \draw[photon] (1.15,1.8) -- (1.65,2.6);
  \node[vertex] at (-1.15,1.8) {}; \node[vertex] at (1.15,1.8) {};
  \draw[cut] (0,-0.25) -- (0,2.65);
  \node[below] at (-1.15,0) {$p$}; \node[below] at (1.15,0) {$p$};
  \node[left,inner sep=2pt] at (-1.65,2.6) {$q$};}
\begin{subfigure}[b]{0.245\textwidth}\centering
\begin{tikzpicture}[fd]
  \draw[quark] (-1.15,0) -- (-1.15,0.75);
  \draw[quark] (-1.15,0.75) -- (-1.15,1.8);
  \draw[quark] (-1.15,1.8) -- (-0.45,1.8);
  \draw[quark=0.65] (-0.45,1.8) -- (1.15,1.8);
  \draw[quark] (1.15,1.8) -- (1.15,0);
  \draw[gluon] (-1.15,0.75) to[out=0,in=-90] (-0.45,1.8);
  \toneDIS
  \node[vertex] at (-1.15,0.75) {}; \node[vertex] at (-0.45,1.8) {};
\end{tikzpicture}
\caption{虚修正}
\end{subfigure}\hfill
\begin{subfigure}[b]{0.245\textwidth}\centering
\begin{tikzpicture}[fd]
  \draw[quark] (-1.15,0) -- (-1.15,0.75);
  \draw[quark] (-1.15,0.75) -- (-1.15,1.8);
  \draw[quark=0.75] (-1.15,1.8) -- (1.15,1.8);
  \draw[quark] (1.15,1.8) -- (1.15,0.75);
  \draw[quark] (1.15,0.75) -- (1.15,0);
  \draw[gluon] (-1.15,0.75) -- (1.15,0.75);
  \toneDIS
  \node[vertex] at (-1.15,0.75) {}; \node[vertex] at (1.15,0.75) {};
\end{tikzpicture}
\caption{初态辐射}
\end{subfigure}\hfill
\begin{subfigure}[b]{0.245\textwidth}\centering
\begin{tikzpicture}[fd]
  \draw[quark] (-1.15,0) -- (-1.15,1.8);
  \draw[quark] (-1.15,1.8) -- (-0.6,1.8);
  \draw[quark=0.27] (-0.6,1.8) -- (0.6,1.8);
  \draw[quark] (0.6,1.8) -- (1.15,1.8);
  \draw[quark] (1.15,1.8) -- (1.15,0);
  \draw[gluon] (-0.6,1.8) arc[start angle=180,end angle=0,radius=0.6];
  \toneDIS
  \node[vertex] at (-0.6,1.8) {}; \node[vertex] at (0.6,1.8) {};
\end{tikzpicture}
\caption{末态辐射}
\end{subfigure}\hfill
\begin{subfigure}[b]{0.245\textwidth}\centering
\begin{tikzpicture}[fd]
  \draw[quark] (-1.15,0) -- (-1.15,0.75);
  \draw[quark] (-1.15,0.75) -- (-1.15,1.8);
  \draw[quark=0.3] (-1.15,1.8) -- (0.6,1.8);
  \draw[quark] (0.6,1.8) -- (1.15,1.8);
  \draw[quark] (1.15,1.8) -- (1.15,0);
  \draw[gluon] (-1.15,0.75) to[out=0,in=-90] (0.6,1.8);
  \toneDIS
  \node[vertex] at (-1.15,0.75) {}; \node[vertex] at (0.6,1.8) {};
\end{tikzpicture}
\caption{干涉项}
\end{subfigure}

\bigskip
\begin{subfigure}[b]{0.42\textwidth}\centering
\begin{tikzpicture}[fd]
  \coordinate (A) at (0,-0.7); \coordinate (B) at (0,0.7);
  \draw[gluon] (-1.5,-0.7) -- (A);
  \draw[photon] (-1.5,0.7) -- (B);
  \draw[quark] (1.5,-0.7) -- (A);
  \draw[quark] (A) -- (B);
  \draw[quark] (B) -- (1.5,0.7);
  \node[vertex] at (A) {}; \node[vertex] at (B) {};
  \path[mom={$q$}] (-1.5,0.7) -- (B);
  \path[mom'={$p$}] (-1.5,-0.7) -- (A);
  \path[mom={$k_1$}] (B) -- (1.5,0.7);
  \path[mom'={$k_2$}] (A) -- (1.5,-0.7);
  \path[mom'={$k_1-q$}] (A) -- (B);
  \node[left] at (-1.5,0.7) {$\gamma^*$};
  \node[left] at (-1.5,-0.7) {$g$};
  \node[right] at (1.5,0.7) {$q$};
  \node[right] at (1.5,-0.7) {$\bar q$};
\end{tikzpicture}
\caption{$\gamma^*g\to q\bar q$，$t$ 道}
\end{subfigure}\hspace{1em}
\begin{subfigure}[b]{0.42\textwidth}\centering
\begin{tikzpicture}[fd]
  \coordinate (A) at (0,-0.7); \coordinate (B) at (0,0.7);
  \draw[gluon] (-1.5,-0.7) -- (A);
  \draw[photon] (-1.5,0.7) -- (B);
  \draw[quark] (1.5,0.7) -- (B);
  \draw[quark] (B) -- (A);
  \draw[quark] (A) -- (1.5,-0.7);
  \node[vertex] at (A) {}; \node[vertex] at (B) {};
  \path[mom={$q$}] (-1.5,0.7) -- (B);
  \path[mom'={$p$}] (-1.5,-0.7) -- (A);
  \path[mom={$k_2$}] (B) -- (1.5,0.7);
  \path[mom'={$k_1$}] (A) -- (1.5,-0.7);
  \path[mom={$q-k_2$}] (B) -- (A);
  \node[left] at (-1.5,0.7) {$\gamma^*$};
  \node[left] at (-1.5,-0.7) {$g$};
  \node[right] at (1.5,0.7) {$\bar q$};
  \node[right] at (1.5,-0.7) {$q$};
\end{tikzpicture}
\caption{$\gamma^*g\to q\bar q$，$u$ 道}
\end{subfigure}
\caption{DIS 在 $\order{\as}$ 的部分子图。(a)--(d) 为 $\gamma^*q$ 道的割图（左侧为振幅，右侧为共轭振幅）：
(a) 一圈顶点修正与树图的干涉（外线自能在维数正规化中无标度，未画出）；(b) 初态辐射的平方；(c) 末态辐射的平方；
(d) 初态辐射与末态辐射的干涉。(a) 与 (d) 各有一个未画出的镜像图（左右互换，即取复共轭），须一并计入。(e)、(f) 为玻色子—胶子聚合 $\gamma^*g\to q\bar q$ 的两个振幅图：虚光子分别与夸克、反夸克耦合，
内线动量分别为 $k_1-q$ 与 $q-k_2$，传播子分母为 $\hat t=(q-k_1)^2$ 与 $\hat u=(q-k_2)^2$；
反夸克线上的箭头为费米子数流向，其动量 $k_2$ 向外。两图之和的平方给出胶子道系数函数。}
\label{fig:dis-oneloop}
\end{figure}
%%% END fig:dis-oneloop
